% ================== Chapitre 2 : conception du service de resolution ========
\chapter{Conception du service de résolution}
\label{ch:conception}

\section{Séparation des rôles autoritatif et récursif}
\label{sec:roles}

Le document 03 impose de distinguer au moins deux rôles de service et précise
qu'une instance récursive partagée non filtrée est insuffisante. Cette
séparation n'est pas une préférence de mise en \oe uvre : elle découle de trois
propriétés incompatibles sur une même instance.

\begin{itemize}
  \item Un serveur autoritatif doit répondre avec autorité pour les zones qu'il
    détient et ne doit rien résoudre d'autre. Un résolveur doit au contraire
    parcourir l'arbre public et conserver un cache.
  \item Les populations de clients diffèrent. L'autoritatif n'a que quatre
    interlocuteurs légitimes : son homologue, les deux résolveurs et l'agent. Le
    résolveur sert au contraire l'ensemble des machines autorisées.
  \item Les surfaces d'attaque diffèrent. Exposer l'autoritatif aux tenants
    reviendrait à leur offrir un point d'interrogation direct de toutes les
    zones hébergées, que seule une politique interne pourrait ensuite
    restreindre.
\end{itemize}

La conception retient donc deux serveurs autoritatifs, \opt{ns1} et \opt{ns2},
qui ne sont joignables que depuis le VNet Services, et deux résolveurs
récursifs, \opt{rec1} et \opt{rec2}, qui constituent le seul point de résolution
offert aux machines clientes. Les résolveurs atteignent l'espace interne par une
redirection explicite de zone, et l'espace public par une récursion ordinaire.

\begin{figure}[H]
\centering
\figArchitecture
\caption{Architecture du service DNS : placement des quatre instances, flux
  autorisés et dépendances communes}
\label{fig:architecture}
\end{figure}

\section{Espace de nommage, zones directes et zones inverses}
\label{sec:nommage}

La zone \zone{gandal.internal} est conservée comme proposition de nommage
interne, conformément au document 03. Le suffixe est pertinent : l'ICANN a
réservé le domaine de premier niveau \zone{internal} à l'usage privé et ne le
délègue pas dans la racine publique. Un nom construit sous ce suffixe ne peut
donc pas entrer en collision avec un nom Internet, ce qui n'est pas le cas d'un
suffixe arbitraire.

\subsection{Découpage retenu}

\begin{itemize}
  \item \zone{gandal.internal} porte le SOA, les enregistrements NS du service
    et les adresses des quatre serveurs, puis délègue ses zones filles.
  \item \zone{infra.gandal.internal} contient l'inventaire d'infrastructure :
    hôtes, orchestrateur, bastion, machines pfSense, services de base et
    collecteurs.
  \item \zone{tenant-0NN.gandal.internal}, pour \opt{NN} de \opt{001} à
    \opt{050}, contient la passerelle du tenant et les noms de ses machines.
\end{itemize}

Les serveurs de noms sont volontairement nommés au niveau de l'apex, sous la
forme \zone{ns1.gandal.internal}, et non à l'intérieur de
\zone{infra.gandal.internal}. Ce choix évite une délégation dont les serveurs
porteraient un nom situé sous la zone déléguée, configuration qui exige des
enregistrements de collage dans la zone parente et complique inutilement la
reconstruction.

Le plan d'adressage du document 02 rend la correspondance entre un tenant et ses
zones entièrement déterministe. En notant $T=50$ le nombre de tenants du
périmètre minimal, elle s'écrit
\begin{align}
  \mathcal{P}(n) &= \texttt{10.30.}(63+n)\texttt{.0/24},
    & 1 &\leqslant n \leqslant T, \label{eq:prefixe}\\[2pt]
  \mathcal{D}(n) &= \texttt{tenant-0}\langle n\rangle\texttt{.gandal.internal},
    \label{eq:directe}\\[2pt]
  \mathcal{R}(n) &= (63+n)\texttt{.30.10.in-addr.arpa},
    \label{eq:inverse}
\end{align}
où $\langle n\rangle$ désigne le numéro du tenant écrit sur trois chiffres. Le
tenant 1 occupe ainsi \ip{10.30.64.0/24} et le tenant 50 occupe
\ip{10.30.113.0/24}, les quatorze préfixes \ip{10.30.114.0/24} à
\ip{10.30.127.0/24} restant en réserve à l'intérieur du \ip{/18}.

\subsection{Résolution inverse}

Toutes les plages du projet appartiennent au bloc \ip{10.30.0.0/16} et sont
découpées en \ip{/24}, en \ip{/28} ou en \ip{/29}. Comme le projet détient
l'intégralité de l'espace, une zone inverse par \ip{/24} suffit à couvrir les
sous-réseaux plus petits qu'elle contient : les quatre \ip{/28} de la plage DMZ
et VPN relèvent tous de \zone{16.30.10.in-addr.arpa}, et les trois \ip{/29} de
synchronisation de \zone{17.30.10.in-addr.arpa}. Le découpage sans classe du
RFC 2317, utile lorsqu'un \ip{/24} est partagé entre plusieurs autorités, n'a
donc aucune raison d'être introduit ici.

\begin{figure}[H]
\centering
\figZones
\caption{Arborescence des zones directes et inverses}
\label{fig:zones}
\end{figure}

Le nombre total de zones à créer se déduit de ce découpage. En comptant la zone
racine interne et la zone d'infrastructure, les onze zones inverses
d'infrastructure, puis une zone directe et une zone inverse par tenant,
\begin{equation}
  N_Z = \underbrace{2}_{\text{apex et infra}}
      + \underbrace{11}_{\text{inverses d'infrastructure}}
      + \underbrace{2T}_{\text{zones des tenants}}
      = 2T + 13 = 113 .
  \label{eq:nzones}
\end{equation}

\begin{longtable}{@{}R{0.21}R{0.30}C{0.13}R{0.36}@{}}
\caption{Inventaire des zones à créer}\label{tab:zones}\\
\hautfilet
\textbf{Catégorie} & \textbf{Nom ou intervalle} & \textbf{Nombre} &
\textbf{Contenu}\\
\medfilet
\endfirsthead
\hautfilet
\textbf{Catégorie} & \textbf{Nom ou intervalle} & \textbf{Nombre} &
\textbf{Contenu}\\
\medfilet
\endhead
\basfilet
\endfoot
Zone racine interne & \zone{gandal.internal} & 1 &
  SOA, NS, adresses des serveurs, délégations\\
Zone d'infrastructure & \zone{infra.gandal.internal} & 1 &
  environ 40 noms, statiques\\
Zones tenant & \zone{tenant-001} à \zone{tenant-050} \zone{.gandal.internal} &
  50 & passerelle statique, machines publiées par l'agent\\
Inverses d'infrastructure & \zone{0, 1, 2, 4, 5, 6, 7, 8, 10, 16, 17
  .30.10.in-addr.arpa} & 11 &
  un \ip{/24} inverse couvre les \ip{/28} et \ip{/29} qu'il contient\\
Inverses tenant & \zone{64} à \zone{113} \zone{.30.10.in-addr.arpa} & 50 &
  tenant $n$ associé à l'octet $63+n$\\
\medfilet
\textbf{Total} & & \textbf{113} &
  volume compatible avec une base locale\\
\end{longtable}

\section{Placement, redondance et domaines de panne}
\label{sec:placement}

L'exigence \exig{net-nfr-001} demande deux instances par service, et le document
Underlay rappelle que la présence de plusieurs serveurs ne garantit pas leur
indépendance si tous partagent le même hyperviseur. Le placement doit donc
satisfaire une contrainte d'anti-affinité par rôle : les deux membres d'une
paire ne résident jamais sur le même hôte.

Avec quatre instances et trois hôtes, un hôte en porte nécessairement deux. La
répartition retenue place \opt{ns1} et \opt{rec2} sur H1, \opt{rec1} sur H2 et
\opt{ns2} sur H3. Elle est préférable à un regroupement sur deux hôtes
seulement, car elle laisse trois instances sur quatre en service après la perte
de H2 ou de H3.

\begin{longtable}{@{}R{0.17}R{0.17}R{0.17}R{0.49}@{}}
\caption{Effet de la perte d'un hôte sur le service de résolution}
\label{tab:pannes}\\
\hautfilet
\textbf{Scénario} & \textbf{Autoritatifs restants} &
\textbf{Récursifs restants} & \textbf{Conséquence}\\
\medfilet
\endfirsthead
\hautfilet
\textbf{Scénario} & \textbf{Autoritatifs restants} &
\textbf{Récursifs restants} & \textbf{Conséquence}\\
\medfilet
\endhead
\basfilet
\endfoot
Fonctionnement nominal & ns1 (H1), ns2 (H3) & rec1 (H2), rec2 (H1) &
  Service complet\\
Perte de H1 & ns2 & rec1 &
  Un serveur de chaque rôle, capacité réduite de moitié ; ns2 devient le seul
  autoritatif et ne reçoit plus de transfert, donc plus de publication\\
Perte de H2 & ns1, ns2 & rec2 & Résolution et publication nominales\\
Perte de H3 & ns1 & rec1, rec2 &
  Publication nominale, réplication suspendue jusqu'au retour de ns2\\
Perte de SW1 & aucune & aucune &
  Hors de portée de la redondance logicielle ; limite matérielle assumée par le
  cadrage\\
\end{longtable}

\begin{encadredecision}
\titreencadredec{Conséquence à retenir pour la reprise}
La perte de H1 supprime le seul serveur autoritatif primaire. La résolution se
poursuit sur \opt{ns2}, qui continue de servir les zones transférées jusqu'à
expiration du paramètre \opt{expire} du SOA, fixé à sept jours. En revanche,
aucune nouvelle publication n'est possible pendant cette période. La procédure
de reprise consiste à restaurer \opt{ns1} depuis l'export des zones, et non à
promouvoir \opt{ns2}, afin de conserver un propriétaire d'écriture unique.

\smallskip
Cette distinction entre continuité du trafic existant et impossibilité
d'opérations nouvelles correspond exactement au résultat attendu du test
\opt{G04} du document 01.
\end{encadredecision}

\section{Politique de durée de vie et cohérence avec le bail DHCP}
\label{sec:ttl}

Les durées de vie ne sont pas un réglage de confort : elles déterminent le délai
pendant lequel une réponse périmée peut encore circuler après la suppression
d'un enregistrement. Deux contraintes les gouvernent. La durée de vie d'un
enregistrement dynamique doit rester nettement inférieure à la durée $B$ du bail
qui l'a produit, et une adresse ne peut être réattribuée qu'après expiration des
caches qui ont pu mémoriser son ancien nom. En notant $Q$ la durée de
quarantaine et $\delta$ une marge de sécurité,
\begin{equation}
  \mathrm{TTL}_A \;\leqslant\; \frac{B}{4}
  \qquad\text{et}\qquad
  Q \;\geqslant\; \mathrm{TTL}_A + \mathrm{TTL}_{\mathrm{nég}} + \delta .
  \label{eq:ttl}
\end{equation}
Avec l'hypothèse $B = 3\,600$ secondes retenue plus bas, le choix
$\mathrm{TTL}_A = \mathrm{TTL}_{\mathrm{nég}} = 300$ secondes et
$\delta = 300$ secondes conduit à $Q = 900$ secondes et respecte les deux
inégalités avec un facteur de sécurité confortable.

\begin{longtable}{@{}R{0.28}R{0.14}R{0.58}@{}}
\caption{Politique de durée de vie}\label{tab:ttl}\\
\hautfilet
\textbf{Paramètre} & \textbf{Valeur} & \textbf{Justification}\\
\medfilet
\endfirsthead
\hautfilet
\textbf{Paramètre} & \textbf{Valeur} & \textbf{Justification}\\
\medfilet
\endhead
\basfilet
\endfoot
SOA \opt{refresh} & 3\,600 s &
  Fréquence de contrôle de ns2 en l'absence de NOTIFY ; la réplication normale
  reste déclenchée par NOTIFY.\\
SOA \opt{retry} & 900 s & Nouvelle tentative après échec de transfert.\\
SOA \opt{expire} & 604\,800 s &
  Durée pendant laquelle ns2 continue de servir une zone sans contact avec ns1 ;
  dimensionnée pour couvrir une panne longue de H1.\\
SOA \opt{minimum} & 300 s &
  Durée de mémorisation d'une réponse négative. Une valeur courte évite qu'un
  nom publié tardivement reste introuvable.\\
Enregistrements d'infrastructure & 3\,600 s &
  Noms stables, modifiés par une opération planifiée.\\
Enregistrements A et PTR tenant & 300 s &
  Publiés et retirés avec le bail, soit $B/12$ sous l'hypothèse
  $B = 3\,600$ s.\\
Quarantaine $Q$ avant réutilisation & 900 s &
  Proposée à l'IPAM. Elle satisfait l'inégalité \eqref{eq:ttl} et répond au
  point laissé \statut{à valider} par le document 03.\\
\end{longtable}

La durée de bail de 3\,600 secondes est une hypothèse de travail : elle relève
de la tâche T004. Si elle est modifiée, les deux inégalités \eqref{eq:ttl}
restent la règle à conserver.

\section{Dimensionnement et contraintes de transport}
\label{sec:dimensionnement}

\subsection{Volume d'enregistrements}

Le volume à servir se déduit du plan d'adressage. Chaque tenant dispose d'un
pool dynamique de 190 adresses, de \ip{.10} à \ip{.199}. En notant $L$ le nombre
de baux simultanés, $T$ le nombre de tenants et $H$ le nombre de noms
d'infrastructure, chaque bail produit un enregistrement A et son PTR, chaque
tenant apporte son SOA, ses deux NS et le couple A et PTR de sa passerelle, et
chaque nom d'infrastructure apporte un A et un PTR :
\begin{equation}
  N_R = 2L + 5T + 2H .
  \label{eq:nrecords}
\end{equation}
La cible historique de 500 baux simultanés, avec $T = 50$ et $H \approx 40$,
conduit à un régime courant
\[
  N_R = 2\times 500 + 5\times 50 + 2\times 40 = 1\,330
\]
enregistrements. Dans l'hypothèse haute où les cinquante pools seraient pleins,
$L_{\max} = 190\,T = 9\,500$ et
\[
  N_R^{\max} = 2\times 9\,500 + 250 + 80 = 19\,330 ,
\]
soit un peu moins de vingt mille enregistrements.

Ce volume ne justifie pas un serveur de base de données dédié. Le moteur de
stockage SQLite intégré à PowerDNS est retenu pour le périmètre minimal, la
réplication étant assurée par transfert de zone et non par duplication de base.
Un moteur relationnel partagé constitue une extension à envisager si la mesure
du débit d'écriture de l'agent le justifie.

\begin{table}[H]
\centering
\small
\caption{Dimensionnement proposé des machines virtuelles}\label{tab:vm}
\begin{tabular}{@{}R{0.14}C{0.13}C{0.12}C{0.12}R{0.49}@{}}
\hautfilet
\textbf{Instance} & \textbf{Processeur} & \textbf{Mémoire} & \textbf{Disque} &
\textbf{Remarque}\\
\medfilet
ns1, ns2 & 2 vCPU & 2 GiB & 20 GiB &
  Base SQLite, journal local, export des zones\\
rec1, rec2 & 2 vCPU & 4 GiB & 20 GiB &
  Mémoire dimensionnée pour le cache d'enregistrements et le cache paquet\\
\basfilet
\end{tabular}
\end{table}

\subsection{Taille des réponses et MTU}

Le document 02 établit la taille utile laissée au paquet IP interne par
l'encapsulation VXLAN sur IPv4, à partir d'une MTU externe de 1\,500 octets :
\begin{equation}
  \mathrm{MTU}_{\text{utile}}
    = 1\,500 - \underbrace{20}_{\text{IPv4}}
             - \underbrace{8}_{\text{UDP}}
             - \underbrace{8}_{\text{VXLAN}}
             - \underbrace{14}_{\text{Ethernet interne}}
    = 1\,450 \ \text{octets}.
  \label{eq:mtu}
\end{equation}
Une réponse DNS qui dépasserait cette taille serait fragmentée ou perdue selon
le comportement des équipements intermédiaires, et ce type de défaut se
manifeste de manière intermittente, donc difficile à diagnostiquer.

La taille annoncée par EDNS est en conséquence fixée à 1\,232 octets sur les
quatre instances. Cette valeur correspond à la recommandation issue du DNS Flag
Day 2020 et laisse une marge
\[
  \mathrm{MTU}_{\text{utile}} - 1\,232 = 218 \ \text{octets}
\]
pour les en-têtes IP et UDP du paquet interne lui-même. Au-delà de ce seuil, la
réponse est tronquée et le client bascule en TCP, comportement prévisible et
observable. Les règles de pare-feu doivent donc autoriser le port 53 en TCP
comme en UDP, point souvent omis.

\section{Profil DNS par tenant et interfaces consommatrices}
\label{sec:profil}

Le modèle IPAM minimal du document 03 comporte un champ \opt{DNS\_profile} dont
le contenu n'est pas spécifié. Ce dossier le définit, car il constitue le point
de passage unique entre la conception DNS, les options distribuées par DHCP et
les paramètres remis aux Templates.

\begin{lstlisting}[caption={Objet profil DNS porté par l'IPAM, instancié pour le
  tenant 7, de préfixe 10.30.70.0/24},label=lst:profil]
{
  "profile_id":         "tenant-standard-v1",
  "tenant_id":          7,
  "resolvers":          ["10.30.4.12", "10.30.4.13"],
  "search_domains":     ["tenant-007.gandal.internal"],
  "forward_zone":       "tenant-007.gandal.internal",
  "reverse_zone":       "70.30.10.in-addr.arpa",
  "record_ttl":         300,
  "publication":        "agent",
  "external_recursion": true,
  "view_tag":           1007
}
\end{lstlisting}

\begin{table}[H]
\centering
\small
\caption{Consommation du profil DNS par les interfaces du lot Services Network}
\label{tab:profil}
\begin{tabular}{@{}R{0.26}R{0.74}@{}}
\hautfilet
\textbf{Interface} & \textbf{Usage du profil}\\
\medfilet
IF-NS-02, vers DHCP &
  \opt{resolvers} alimente l'option 6, \opt{search\_domains} l'option 15 et
  l'option 119. Le serveur Kea ne redéfinit aucune de ces valeurs localement.\\
IF-NS-06, vers Templates &
  \opt{resolvers}, \opt{search\_domains} et la MTU utile sont remis au mécanisme
  d'initialisation de la machine. Le profil ne contient aucun secret.\\
IF-NS-03, vers VNet et SDN &
  \opt{view\_tag} matérialise le profil de résolution et la liste de contrôle
  d'accès associée au tenant.\\
IF-NS-05, vers Backend par l'agent &
  \opt{forward\_zone} et \opt{reverse\_zone} désignent les deux zones à modifier
  lors d'une publication ou d'un retrait.\\
\basfilet
\end{tabular}
\end{table}

Un second profil, \opt{admin-v1}, désigne les mêmes résolveurs avec le domaine
de recherche \zone{infra.gandal.internal} et l'étiquette de vue 1. Il s'applique
aux machines du plan de gestion et au pool VPN administrateur.
